\documentclass[crop=false]{standalone}
% --- ПРЕАМБУЛА ДЛЯ ПОДДЕРЖКИ РУССКОГО ЯЗЫКА И МАТЕМАТИКИ ---
\usepackage[T2A]{fontenc}
\usepackage[utf8]{inputenc}
\usepackage[russian]{babel}
\usepackage{amsmath}
\usepackage{geometry}
\usepackage{amssymb}
\usepackage{graphicx}
\usepackage{float}
\usepackage{import}
\geometry{a4paper, top=2cm, bottom=2cm, left=2cm, right=2cm}
\newcommand{\sidecomment}[1]{\marginpar{\raggedright\scriptsize\itshape #1}}
\newcommand{\iu}{\mathrm{i}}

\newcommand{\Def}{%
    \text{Def}%
    \hspace{-0.85em}% Сдвиг назад ровно на середину (подберите под шрифт)
    \raisebox{0.1ex}{/}% Черта
    \hspace{1em}% Возврат к концу слова + небольшой отступ
}

\renewcommand{\Re}{\operatorname{Re}}
\renewcommand{\Im}{\operatorname{Im}}

\ifstandalone
    \newcommand{\lecturebreak}{} % В отдельном файле лекции разрыв не нужен
\else
    \newcommand{\lecturebreak}{\clearpage} % В полной книге - начинаем с новой страницы
\fi

\newcounter{lecturenumber} % Создаем счетчик для нумерации лекций
\newcommand{\lecture}[2]{
    \lecturebreak
    \refstepcounter{lecturenumber} % Увеличиваем счетчик на 1
    \par\vspace{2em}\noindent % Отступ сверху и убираем абзацный отступ
    {\Large\bfseries % Делаем заголовок крупным и жирным
    Лекция \thelecturenumber: #1 % Выводим "Лекция N: Тема"
    \hfill % Растягиваем пробел, чтобы дата была справа
    \large\textit{#2}} % Выводим дату курсивом
    \par\vspace{1em} % Небольшой отступ снизу
    \addcontentsline{toc}{section}{Лекция \thelecturenumber: #1} % Добавляем в оглавление
    \par
}

\newcommand{\TwoCols}[2]{%
  \par\noindent % Начать с новой строки без отступа
  \begin{minipage}[t]{0.48\textwidth}
    #1
  \end{minipage}% <-- Важный '%' для удаления пробела
  \hfill % Растягивающийся пробел между колонками
  \begin{minipage}[t]{0.48\textwidth}
    #2
  \end{minipage}% <-- Важный '%' для удаления пробела
  \par%
}

\pagestyle{plain} % Включаем нумерацию страниц\
\begin{document}
\lecture{}{10.03}
Прямая задача — найти движение по известным силам

Обратная задача — найти силы по известному движению
\section{Дополнительные аксиомы}
\begin{enumerate}
    \item \[m\bar w = \bar F \]
    \item \[\bar F_{k m} = -\bar F_{m k} \]
\end{enumerate}
\begin{enumerate}
    \item 
    \Def Сила — вектор, который сопоставляется точке
    
    \item
    \Def $\bar m_O (\bar F) = (\bar r \times \bar F)$ — момент силы
    \[\bar m_A (\bar F) = (\bar r - \bar r_A) \times \bar F\]

    \item Энергия $\backslash$ Работа 
    
    Элементарная работа:
\[d' A = (\bar F , d \bar r) \qquad \Delta A = \bar F \Delta \bar r \qquad \delta A = \bar F \delta \bar r\]
действительное перемещние , возможное перемещние, виртуальное перемещение

    \item
    $N, \bar r_k$ — система точек
    \Def Главный вектор
    \[\bar R = \sum_{k=1}^N \bar F_k\]
    Не всегда равен равнодействующему

    \item 
    Силы  % сделать схемку с двумя стрелками
    \begin{itemize}
        \item Внутренние
        \item Внешние
    \end{itemize}
    \Def Замкнутая система — система, в которой отсутсвует внешние силы
    \Def Незамкнутая система — не замкнутая система
    \[\bar R = \sum_{k=1}^N \bar F_k = \sum_{k=1}^N \bar F_k^{(e)} + \sum_{k=1}^N \bar F_k^{(i)} = \sum_{k=1}^N F_k^{(e)}\]
    Сумма внутренних сил равна нулю, так как точки действуют друг на друга, а действующая равна противодействующей 

    \item 
    \Def Главный момент сил
    \[\bar M_O = \sum_{k=1}^N \bar m_O(\bar F_k) = \sum_{k=1}^{N} (\bar r_k \times \bar F_k)\]
    \[\bar M_A = \sum_{k=1}^N \bar m_A(\bar F_k) = \sum_{k=1}^{N} ((\bar r_k - \bar r_A)\times \bar F_k)\qquad \bar M_B = \sum_{k=1}^N \bar m_B(\bar F_k) = \sum_{k=1}^{N} ((\bar r_k - \bar r_B)\times \bar F_k)\]
    \[\bar M_A - \bar M_B = \sum_{k=1}^{N} (\bar r_B - \bar r_A) \times \bar F_k = (\bar r_B - \bar r_A) \times \bar R\]
    Правило переноса полюса главного момента сил системы
    \[\boxed{\bar M_A = \bar M_B + (\bar r_B - \bar r_A) \times \bar R} = \bar M_B + \bar R \times \overline{BA}\]
    
    \keypoint
    \[\bar M_O = \sum_{k=1}^{N} m_O(\bar F_k) = \sum_{k=1}^{N} (\bar r_k \times \bar F_k) = \sum_{k=1}^{N} \bar r_k \times \bar F_k^{(e)} \]
    Доказательство, что момент внутренних сил нуль:
    \[\sum_{k=1}^{N} \bar r_k \times \left( \sum_{m=1}^{N} \bar F_{k m}^{(i)} \right) = 
    \sum_{k, m}^{N} \bar r_k \times \bar F_{k m}^{(i)} 
    = \sum_{k < m} \bar r_k \times \bar F_{k m}^{(i)} + \sum_{k > m} \bar r_k \times \bar F_{k m}^{(i)} 
    = \sum_{k < m} \bar r_k \times \bar F_{k m}^{(i)} + \sum_{m > k} \bar r_k \times -\bar F_{m k}^{(i)} = 0
    \] % fixme дописать
\end{enumerate}
\section{Общие уравнения динамки}

% fixme шарик на нитке фото в телефоне
\[\begin{cases}
    \ddot{\bar r} m = \bar w m = \bar F + \bar R_0, \\
    x^2 + y^2 - L^2 = 0
\end{cases} \qquad \begin{cases}
    \ddot{\bar r} m = \bar w m = \bar F + \bar R_O\\
    \dot x^2 + x \ddot x + \dot y^2 + y \ddot y = 0
\end{cases}\]
здесь $\bar R_0$ — сила реакции связи

Силы:
\begin{itemize}
    \item Активные
    \item Реакции связи
\end{itemize}

\[k = 1, \dots, N \qquad \bar w_k m_k = \bar F_k + \bar R_k\]

\Def Свзяи называются идеальными, если их элементарная работа совершённая на(любых) виртуальных перемещениях равна нулю
\[\delta A =\sum_{k=1}^{N}\bar R_k \delta r_K = 0\]

\[\bar w_k m_k - \bar F_k = \bar R_k \qquad \sum_{k=1}^{N} (\bar w_k m_k - \bar F_k) \delta \bar r_k = \sum_{k=1}^{N} \bar R_k \delta \bar r_k = 0\]
Общее уравнение динамики: 
\[\boxed{\sum_{k=1}^{N} (\bar w_k m_k - \bar F_k) \delta \bar r_k = 0} \qquad \text{где }\bar F_k\text{ — активные силы}\]

\[\sum_{k=1}^{N} (\bar F_k - \bar w_k m_k) \delta \bar r_k = 0 \qquad -\bar w_k m_k = J_k \qquad J_k \text{ — сила инерции}\]

Движение системы таково, что совокупная работа активных сил и сил инерции на виртуальных перемещениях равна нулю

\noindent \keypoint В неинерциальной системе отсчёта
\[
\bar w_k = \bar w_k^{\text{отн}} + \bar w_k^{\text{кор}} + \bar w_k^{\text{пер}}\]
Тогда 
\[\sum_{k=1}^{N} \left( \bar F_k - \underbracket{\bar w_k^{\text{кор}} m_k}_{J_k^{\text{кор}}} - \underbracket{\bar w_k^{\text{пер}} m_k}_{J_k^{\text{пер}}} - \underbracket{\bar w_k^{\text{отн}} m_k}_{J_k^\text{отн}} \right) \delta \bar r_k = 0\]
\keypoint Любая связь может быть заменена её реакцией
\[\sum_{k=1}^{N} \left(\bar F_l + \bar G_k - \bar w_k m_k \right) \delta \bar r_k =0 \qquad \text{где }\bar G_k \text{ — реакция неидеальных связей}\]

\section{Равновесие}
\Def Положение равновесия механической системы — положение системы, в котором она будет неподвижна на протяжении всего времени, если она изначально была неподвижна


\noindent \keypoint

Система нафходится в равновесии 

$\iff$

Работа активных сил на виртуальных перемещениях равна нулю
\[\boxed{\sum_{k=1}^{N} \bar F_k \delta r_k = 0} \qquad \text{Принцип виртуальных перемещений}\]
\begin{proof}
    $\bar v_k = 0 \quad \bar w_k = 0$ и общее уравнение динамики
\end{proof}

\noindent \keypoint

для голономной системы

Принцип виртуальных перемещений: все обобщённые силы равны нулю
\begin{proof}
    \[
        \delta \bar r_k = \sum_{k=1}^{n} \frac{\partial \bar r_k}{\partial q_i
        } \delta q_i 
        \qquad
        \sum_{k=1}^{N} \bar F_k \left(\sum_{i=1}^{n} \frac{\partial \bar r_k}{\partial q_i
        } \delta q_i \right) = \sum_{i=1}^{n} \left(\underbracket{\sum_{k=1}^{N} \bar F_k \frac{\partial \bar r_k}{\partial q_i}}_{Q_i}\right) \delta q_i \qquad Q_i \text{ — обобщенная сила (Опр)}\]
        $\delta q_i$ независимы в голономной системе:
        \[\left(\begin{cases}
                \delta q_k \neq 0 \\
                \delta q_i = 0 \; i \neq k
            \end{cases}
            \implies Q_k = 0\right) \implies \left(\sum_{k=1}^{N} \bar F_k \delta r_k = \sum_{i=1}^{n} Q_i \delta q_i =0 \iff Q_i = 0\right)\]
\end{proof}
            

\noindent \keypoint Размерность Обобщённых сил
$\left[ Q_i\right] = \frac{\text{Дж}}{\left[q_i \right]}$
\end{document}